\scanpage{137}
\textcircled{b} $g\leq1$, i.e. $g=0,1$. Est-ce que
\[
H^1(\widehat X,\underline\lambda)\longrightarrow H^1(\widehat X,\underline\lambda(D))
\quad\text{est inj. ?}
\]
ou encore

\diagram{\begin{tikzcd}[ampersand replacement=\&,column sep=large,row sep=small]
H^0(\widehat X,\underline\omega^{-1})\&H^0(\widehat X,\underline\omega^{-1}(-D))\arrow[l]\\
\underline{\operatorname{Lie}}\,\underline{\operatorname{Aut}}(\widehat X)\arrow[u,equal]\&
\underline{\operatorname{Lie}}\,\underline{\operatorname{Aut}}(X)\arrow[u,equal]\\
{\underline t_{\widehat X/\mathbb C}\quad\text{fibré tangent}}\arrow[uu,equal]\&
{\underline t_{\widehat X/\mathbb C}(-D),\quad\text{champs \unsure{tangts|tangents} s'annulant sur }D}\arrow[uu,equal]
\end{tikzcd}}

\deleted{injectif} surjectif ? Si son conoyau \unsure{$\rightarrow$}\inserted{\unsure{sans flèche|sans plus|sans pôles}}
$Q(X)$, on aura
\[
0\longrightarrow H^0(\widehat X,\underline\omega^2)\longrightarrow
H^0(X,\underline\omega^2)^{\mathrm{rég}}\longrightarrow
H^0(D,\underline\omega|D)\longrightarrow\check Q(X)\longrightarrow0
\]

\uline{Cas} $g=1$, on \deleted{\unsure{dim}} $\operatorname{rg} H^0(\widehat X,\underline t_{\widehat X/\mathbb C})=\deleted{\unsure{\operatorname{rg}H^1\mathbb C}}1$
et si $D\ne0$ i.e. $X\ne\widehat X$ i.e. $X$ non complète,
on a $Q(X)\simeq H^0(\widehat X,\underline\omega^{-1})\simeq t_{\widehat X}$
(\deleted{\unsure{espace tangent}}\inserted{Algèbre de Lie} de
la Jacobienne) d'où $\check Q(X)=\omega_{\widehat X}=\check t_{\widehat X}$ (formes
diff. inv. sur jacobienne). \unsure{Interprét.|Interprétation} la plus \unsure{nette} :
on regarde les Pb modulaires des \uline{courbes}
\uline{elliptiques} (à \uline{pt marqué}) avec $\nu$ pts
distincts \uline{de somme nulle} \deleted{\unsure{et celà est que si elles sont}}\inserted{(ou ce qui revient au même,}
courbes elliptiques \unsure{princ.} homogènes et
$\nu$ pts \deleted{marqués}\inserted{(non ordonnés)} distincts et un pt choisi
de \unsure{tel} que ce soit la somme desdits \ldots.

\uline{Cas} $g=0$ \inserted{Alors $H^0(\widehat X,\underline t_{\widehat X/\mathbb C})=0$,}
$H^0(\widehat X,\underline\omega^{-1}_{\widehat X/\mathbb C})$ est de dim $3$, et
suivant les cas

\diagram{$\begin{array}{|ll}
\nu=1&\dim H^0(\widehat X,\underline t_{\widehat X/\mathbb C}(-D))=2\\
\nu=2&\rule{3.7cm}{0.35pt}=1\\
\nu\geq3&\rule{3.7cm}{0.35pt}=0
\end{array}$}

donc $Q(X)$ est respectiv$^{\mathrm{t}}$ de dim $1,2$ ou $3$.
Donc pour $\uline{0\leq\nu\leq3}$, on aura $\dim\check Q(X)=\operatorname{card}\operatorname{supp}D
=\dim H^0(D,\underline\omega|D)$
% archive sheet number: 136
